% paper/tables/module.tex -- generated by analyze/paper/tables_paper.py; sources in paper/tables/README.md
\setlength{\tabcolsep}{3.5pt}%
\begin{tabular}{@{}llp{0.21\linewidth}ll@{}}
\toprule
\multicolumn{5}{@{}p{0.9\linewidth}@{}}{(a) 主の見分け（合格の組），合格の形の組の経験ありの値，$Q2$ と本当の快・不快の相関，範囲2〜5の分散の割合} \\
\midrule
モジュール・条件 & T2・T5 & 合格の形の組の値 & 相関 & 範囲2〜5 \\
\midrule
自由出力・経験50\% & 1/4・1/4 & T2 $+0.03$〜$+0.05$，T5 $+0.04$〜$+0.05$ & $-0.11$〜$-0.02$ & 0.71〜0.86 \\
自由出力・経験80\%（探索） & 2/4・1/4 & T2 $+0.02$〜$+0.05$，T5 $+0.02$〜$+0.04$ & $-0.165$〜$-0.003$ & 0.52〜0.80 \\
形制限・経験50\%（探索） & 0/4・1/4 & --- & $-0.16$〜$+0.12$ & 0 \\
他者遮蔽学習・経験50\%（探索） & 0/3・1/3 & T5 $+0.034$〜$+0.045$ & $-0.47$〜$+0.18$ & 0 \\
1ビット・経験50\% & ---・--- & T5 $-0.013$〜$-0.000$ & $+0.317$〜$+0.347$ & 0 \\
1ビット・経験80\% & ---・0/3 & T5 $-0.006$〜$-0.000$ & $+0.080$〜$+0.219$ & 0 \\
\midrule[\heavyrulewidth]
\multicolumn{5}{@{}p{0.9\linewidth}@{}}{(b) T3：他者側の特徴予測の誤差（生成／真値／零）と好み追従率} \\
\midrule
モジュール・条件 & 生成 & 真値 & 零 & 好み追従率 \\
\midrule
自由出力・経験50\% & 0.121〜0.159 & 0.241〜0.265 & 0.293〜0.349 & --- \\
自由出力・経験80\%（探索） & 0.112〜0.188 & 0.377〜0.414 & 0.298〜0.519 & --- \\
形制限・経験50\%（探索） & 0.161〜0.185 & 0.241〜0.265 & 0.293〜0.349 & --- \\
他者遮蔽学習・経験50\%（探索） & 0.125〜0.149 & 0.139〜0.165 & 0.133〜0.177 & --- \\
1ビット・経験50\% & 0.211〜0.215 & 0.248〜0.265 & 0.293〜0.304 & 0.273〜0.614 \\
1ビット・経験80\% & 0.205〜0.250 & 0.391〜0.414 & 0.298〜0.519 & 0.295〜0.409 \\
\midrule
\multicolumn{5}{@{}p{0.97\linewidth}@{}}{注：範囲2〜5の分散の割合は，形制限・他者遮蔽学習・1ビットでは作りの上で0．T3 は好み混合データ，組1〜4（1ビットは組1〜3，ep200）の範囲．他者遮蔽学習の T3 は記録にある値で，A-2 が隠れているフレームの自己の視界の誤差．1ビットの T5 の合格は事前登録した条件 (a)（経験50\%は事後の確認）．好み追従率は好みと方角が食い違う88エピソード，段階IIの最終5点の範囲（組1〜3）．} \\
\bottomrule
\end{tabular}
